\chapter{Performer 与随机特征注意力}

\section{概述}

Performer 由 Choromanski 等人于 2020 年提出，是另一种重要的线性注意力方法。与 Linear Transformer 使用 elu+1 映射不同，Performer 提出了 FAVOR+ 机制（Fast Attention Via positive Orthogonal Random features），通过随机特征近似 softmax 核函数。

\section{FAVOR 机制}

FAVOR 的核心思想是使用随机傅里叶特征（Random Fourier Features, RFF）来近似 softmax 核函数。

\subsection{softmax 核的随机特征近似}

考虑 softmax 注意力中的核函数：

\begin{equation}
    \kappa(\mathbf{q}, \mathbf{k}) = \exp(\mathbf{q}^\top \mathbf{k})
\end{equation}

这里省略了缩放因子和归一化。对于查询和键向量都经过适当归一化后，有：

\begin{equation}
    \|\mathbf{q} - \mathbf{k}\|^2 = \|\mathbf{q}\|^2 + \|\mathbf{k}\|^2 - 2\mathbf{q}^\top \mathbf{k}
\end{equation}

因此：

\begin{equation}
    \exp(\mathbf{q}^\top \mathbf{k}) \propto \exp\left(-\frac{1}{2}\|\mathbf{q} - \mathbf{k}\|^2\right)
\end{equation}

这正是高斯核（RBF 核）的形式。根据 Bochner 定理，高斯核可以通过随机傅里叶特征近似：

\begin{equation}
    \exp(\mathbf{q}^\top \mathbf{k}) \approx \mathbb{E}_{\boldsymbol{\omega}}\left[\exp(\boldsymbol{\omega}^\top \mathbf{q}) \exp(\boldsymbol{\omega}^\top \mathbf{k})\right]
\end{equation}

其中 $\boldsymbol{\omega} \sim \mathcal{N}(\mathbf{0}, \mathbf{I})$。

\subsection{正交随机特征}

FAVOR+ 的关键改进是使用正交随机特征来降低近似方差。具体地：

\begin{equation}
    \phi(\mathbf{x}) = \frac{1}{\sqrt{m}} \begin{bmatrix}
        \exp(\boldsymbol{\omega}_1^\top \mathbf{x}) \\
        \vdots \\
        \exp(\boldsymbol{\omega}_m^\top \mathbf{x})
    \end{bmatrix}
\end{equation}

其中 $\boldsymbol{\omega}_i$ 通过正交化过程采样，使得它们相互正交。正交特征相比独立随机采样具有更低的方差，从而可以用更少的特征维度 $m$ 达到相同的近似精度。

\section{FAVOR+ 的具体实现}

FAVOR+ 使用以下特征映射：

\begin{equation}
    \phi_i(\mathbf{x}) = \frac{1}{\sqrt{m}} \exp\left(-\frac{\|\mathbf{x}\|^2}{2}\right) \exp(\boldsymbol{\omega}_i^\top \mathbf{x}), \quad i = 1, \ldots, m
\end{equation}

其中 $\boldsymbol{\omega}_i$ 通过以下方式生成：

\begin{enumerate}
    \item 从标准正态分布采样 $m$ 个向量 $\mathbf{r}_1, \ldots, \mathbf{r}_m \sim \mathcal{N}(\mathbf{0}, \mathbf{I})$
    \item 对它们进行 Gram-Schmidt 正交化，得到 $\boldsymbol{\omega}_1, \ldots, \boldsymbol{\omega}_m$
\end{enumerate}

\section{理论保证}

Performer 论文证明了以下理论结果：

\begin{itemize}
    \item FAVOR+ 的特征映射以任意精度近似 softmax 核，所需特征数 $m = \mathcal{O}(\log N)$
    \item 正交随机特征的方差比独立随机特征低，近似误差的方差上界为 $\mathcal{O}(1/m)$
    \item 对于有界输入，FAVOR+ 注意力与标准注意力的输出差异可以被严格控制
\end{itemize}

\section{与其他线性注意力的比较}

\begin{table}[H]
    \centering
    \begin{tabular}{l|ccc}
        \toprule
        方法 & 特征映射 & 特征数 & 归一化 \\
        \midrule
        Linear Transformer & $\text{elu}(\mathbf{x}) + 1$ & $d$ & 显式 \\
        Performer (FAVOR+) & 随机指数特征 & $m > d$ & 隐式/显式 \\
        \bottomrule
    \end{tabular}
    \caption{线性注意力方法比较}
\end{table}

Performer 的 FAVOR+ 相比 Linear Transformer 的 elu+1 映射：

\begin{itemize}
    \item 优点：理论上更接近标准 softmax 注意力的行为
    \item 缺点：特征维度 $m$ 通常大于 $d$，计算开销稍大
\end{itemize}

\section{训练技巧}

Performer 在训练中使用了以下技巧来提高稳定性：

\begin{enumerate}
    \item 重参数化：在特征映射前对输入进行归一化
    \item 广义注意力框架：允许特征映射输出负值，通过额外的归一化处理
    \item 混合训练：先用标准注意力预训练，再替换为线性注意力微调
\end{enumerate}
